\clearpage
\section{The Propers: Themes and Movement}
\zlabel{triptych:concise:themes:start}
God's gifts establish real duties without making their recipients masters of the gift. Cyrus receives power, Caesar receives tribute, and the Thessalonians receive a Gospel that becomes active faith. The worshipper must both distinguish what is owed and become willing to render it. One question concerns the rightful place of earthly goods; the other concerns the transformation of the person who uses them. The Mass's movement from petition to reception brings both questions into worship.

\subsection{Shelter and strength belong to the giver}
The Entrance begins with need. Its Psalm asks to be heard and sheltered; its speaker does not approach God as someone whose strength already secures his life. Augustine understands the voice as Christ with his body, the Church, and its confident cry as the fruit of earlier divine hearing. Confidence thus grows from mercy. The Collect carries dependence into the will itself: sincere service requires more than knowing what an obligation is. It requires a disposition for which the assembly asks.

Isaiah makes the distinction between receiving and recognizing startlingly concrete. Cyrus is named and strengthened for Israel's sake, yet does not know the Lord who equips him. Jerome insists that the Persian king remains the literal addressee: ignorance cannot be transferred to Christ. Nevertheless, Cyrus's real work of liberation can foreshadow the true liberator. The king's usefulness does not make him divine. Aquinas presses the further difference between benefit for the people and recognition in the ruler himself. Successful action does not establish the ruler's conversion.

\subsection{Glory returns to God; duties remain real}
The responsorial Psalm calls all peoples to honour the Creator. The nations' gods are not his equals, and their rulers do not delimit his rule. Augustine hears a new love singing through a worldwide house. His warning reaches the preacher: to announce God's glory while taking the credit oneself is to reverse the song's direction. The temptation to appropriate a received power appears inside religious service as well as outside it.

Jesus' tribute answer likewise orders a genuine earthly claim beneath God's claim. Chrysostom explains ``render'' as repayment of dues, while excluding demands that injure godliness. Aquinas then distinguishes the coin bearing Caesar's image from the soul bearing God's image. His mystical reading reaches the whole person, without dissolving the literal debt. The person who gives Caesar a coin does not thereby acquire a part of life exempt from God; the person who worships God does not thereby escape every civic obligation.

This is the first interpretation's answer: created goods can be received, used and returned without becoming ultimate goods. Its moral fruit is truthful allegiance. It needs the second interpretation's question, however, because an accurate account of duties can coexist with a divided will. The Gospel's questioners say true things about Jesus while trying to trap him. Right words and right service have come apart.
\clearpage
\subsection{The received word becomes a visible life}
Paul thanks God for people who act. Their faith works, their love labours and their hope endures. Chrysostom treats those words as costly realities: their fidelity persists under pressure, rather than waiting for every difficulty to disappear. Aquinas explains that these fruits manifest election without purchasing it. Thanksgiving therefore recognizes human effort while refusing to make that effort its own origin.

The acclamation joins shining in the world to the word of life. Aquinas describes the saints' light as participation in the Word's light. In the surrounding passage, he explains that divine action reaches willing as well as doing, without abolishing freedom. Chrysostom accents a life that serves willingly instead of murmuring. The second interpretation's answer is therefore active dependence: receiving grace makes a person's response possible, and that response becomes visible in service. Capacity without acknowledgment, exemplified by Cyrus, is not the same as grateful fidelity.

At the altar, the Prayer over the Offerings places purification beside service. Having gifts to offer does not mean that the giver has no further need. Mixed motives can enter an outwardly good action; offering becomes a place where the servant receives help. The prayer belongs with the Collect's concern for willing and the final prayer's concern for growth. Together they locate service within a life still being changed.

\subsection{Mercy feeds; the Son gives himself}
The Communion alternatives bring distinct images. The Psalm option places hope in the Lord's merciful care. Its wider context contrasts that care with armies and horses that cannot finally save. Augustine interprets famine through the need for truth; Aquinas distinguishes bodily, spiritual and sacramental nourishment. Fear joined to hope keeps the recipient from both presumption and despair. Neither interpretation promises freedom from bodily suffering.

The Mark option selects Christ's gift of his life as ransom. His teaching on service belongs to the surrounding Gospel context, not additional appointed antiphon wording. Bede joins the blood shed for the world to willing belief in his explanation of ``for many.'' The gift comes first; believing reception and the imitation of Christ's humility follow. Chrysostom's account of Matthew's parallel explains how the Lord's abasement overturns the disciples' ambition. Service is a response to redemption, not another redemption performed by the disciple.

The Prayer after Communion directs present participation toward eternal gifts. For ordered allegiance, this means earthly goods are freed from the impossible task of saving finally. For willing service, it means growth remains real while fulfilment remains ahead. These are complementary consequences: neither a correct allocation of duties nor a sincere beginning makes the pilgrim self-sufficient. The whole person belongs to God, and the whole person still needs the help that brings that belonging to fruition.
\zlabel{triptych:concise:themes:end}
